\documentclass[12pt,a4paper]{article}
\usepackage[utf8]{inputenc}
\usepackage{amsmath}
\usepackage{amsfonts}
\usepackage{amssymb}
\usepackage{geometry}
\usepackage{booktabs}
\usepackage{hyperref}

\geometry{margin=1in}

\begin{document}

\begin{center}
    \large \textbf{The Principle of Planck’s Field: Resolving the Dark Matter Paradox via Density Satiation and the Metric Impedance Scalar} \\
    \vspace{1em}
    \normalsize \textbf{Robin Skavberg} \\
    \normalsize Independent Researcher \\
    \normalsize ORCID: 0009-0003-9126-6196 \\
    \normalsize Email: skavberg@gmail.com \\
    \vspace{1em}
    \normalsize February 2026
\end{center}

\begin{abstract}

This paper identifies Dark Matter as a state of maximum density satiation within the Total Field. We demonstrate that any mass, regardless of scale, transitions from a baryonic (vibrational) state to a dark (non-vibrational) state upon reaching its satiation radius $R_{sat}$. At this limit, the Metric Impedance Scalar $\zeta$ becomes infinite, resulting in a "Frozen Core" characterized by a total absence of internal degrees of freedom. We provide the mathematical proof that the neutrino’s known mass-to-radius ratio places it naturally within this satiated regime, effectively classifying the neutrino as the fundamental subatomic constituent of the Dark Matter energy density.
\end{abstract}

\section{Background}
In traditional physics, baryonic matter is defined as any matter composed primarily of baryons (such as protons and neutrons), making up the stars, planets, and everything we see. However, within this framework, it is defined by its mechanical relationship with the vacuum. Unlike dark matter, which is matter at maximum density, baryonic matter exists at densities below the satiation threshold ($\rho_{sat}$). Because it has not displaced the field to its absolute limit, the field retains its ``slack'' or elasticity.

\subsection{Vibrational Degrees of Freedom}
Because the field is not fully displaced, it can support wave-like perturbations. These vibrations manifest as:
\begin{itemize}
    \item \textbf{Thermal Energy:} The movement and oscillation of particles within the field.
    \item \textbf{Electromagnetic Interaction:} The ability to emit, absorb, and reflect light.
\end{itemize}

\subsection{Comparison of States}
The difference between baryonic and dark matter is essentially a difference in field impedance:

\begin{center}
\begin{tabular}{lll}
\toprule
\textbf{Feature} & \textbf{Baryonic Matter} & \textbf{Dark Matter (Frozen Cores)} \\
\midrule
Density & Low to Medium ($\rho_m < \rho_{sat}$) & Maximum ($\rho_m = \rho_{sat}$) \\
Field Displacement & Partial & Maximum \\
Vibrations & Active (supports heat/light) & Zero (Absolute Zero) \\
G-Coupling & High/Standard & Minimum (Suppressed) \\
\bottomrule
\end{tabular}
\end{center}

In short, baryonic matter is matter that is still ``in conversation'' with the field's vibrational mechanics, whereas dark matter is matter that has been compressed by the Planck Field into a state of silent, rigid satiation.

\section{Dark Matter as Maximum Density Matter}
Within the \textit{Principle of Planck's Field}, both baryonic matter and dark matter are constituents of the field. The distinction lies in the density of the matter and the resulting displacement of the field. Dark matter is defined as matter that has reached the threshold of maximum density.

The gravitational constant $G$ is the dynamic resultant of the field's restorative stiffness ($c^2$) and the matter-energy density ($\rho_m$):
\begin{equation}
G(\rho_m) = \frac{c^2}{8\pi\rho_m}
\end{equation}
When matter reaches the satiation threshold $\rho_{sat}$, it displaces the field to its maximum allowable value. In this state, the matter is ``frozen'' by the field's restorative limit, creating a core that is structurally rigid and incapable of internal oscillation.

\section{Thermodynamic Satiation and Absolute Zero}
The mechanical state of matter at maximum density dictates its thermodynamic properties. Because the matter has displaced the field to its limit, there is no capacity for wave-like perturbations or internal vibrations.
\begin{itemize}
    \item \textbf{Vibrational Satiation:} In Planck's Field, thermal energy is the manifestation of field vibrations. 
    \item \textbf{Absolute Zero:} Matter that has reached maximum density cannot support these vibrations due to the satiation of the field. This lack of vibration is synonymous with the thermodynamic value of absolute zero.
\end{itemize}
This physical state accounts for the ``dark'' profile observed in astrophysics; without vibrational degrees of freedom, the matter cannot emit or absorb radiation, remaining electromagnetically silent while maintaining its gravitational influence.

\section{Mathematical Derivation of Scale-Invariance}
We demonstrate that the transition of matter into this frozen, maximum-density state is a function of density and field displacement, independent of total mass.

In the context of General Relativity and the Principle of Planck's Field, the Schwarzschild boundary condition (or Schwarzschild radius) defines the physical limit where a mass becomes so compressed that the escape velocity equals the speed of light.

The Schwarzschild boundary condition is framed as a radius because General Relativity assumes $G$ is a fixed constant. If $G$ never changes, then the only way to reach the "point of no return" is to cram mass into a specific, tiny volume. However, in the Principle of Planck's Field, the Schwarzschild condition is revealed to be a density limit ($10^{25}$ kg/m$^3$). As The Planck Field Theory correlates the density of mass wherein gravity is the field response to matter in relation to the speed of light, we are no longer constrained by a fixed radius of matter. 

Here is the mechanical breakdown of why the ``radius'' is actually a mask for the field's ``density threshold'':

\subsection{The Variable $G$ Correction}
In the Planck Field framework, $G$ is not a universal constant; it is a resultant of the field's response to matter:
\begin{equation}
G = \frac{c^2}{8\pi\rho_m}
\end{equation}
By adding this definition of $G$ back into the classical Schwarzschild radius formula ($R_s = 2GM/c^2$), something remarkable happens:
\begin{equation}
R_s = \frac{2 \left(\frac{c^2}{8\pi\rho_m}\right) M}{c^2}
\end{equation}
The $c^2$ terms cancel out, and we are left with a relationship between radius, mass, and density:
\begin{equation}
R_s = \frac{M}{4\pi\rho_m}
\end{equation}

\subsection{The Density Satiation Point}
Since $M = \text{Density} \times \text{Volume}$ (and Volume for a sphere is $\frac{4}{3}\pi R^3$), we substitute that into the equation:
\begin{equation}
R = \frac{\rho_m \cdot \frac{4}{3}\pi R^3}{4\pi\rho_m}
\end{equation}
Everything cancels out except for the geometric constants. This proves that the Schwarzschild boundary occurs when the density $\rho_m$ reaches the field's maximum displacement limit.

The Schwarzschild boundary is not just a point of no return for light; it is the geometric threshold where matter reaches maximum density at the satiation limit of the field.

\textbf{Physical Impossibility of the Singularity:}\\
The revelation of the Schwarzschild boundary as a density-driven satiation limit eliminates the mathematical necessity of the gravitational singularity. In classical models, the lack of a counter-force leads to an infinite collapse; however, within the Principle of Planck's Field, the restorative pressure of the vacuum at the satiation threshold acts as a fundamental structural limit. Once matter reaches the density $\rho_{sat}$, the field transitions into a rigid, non-vibrational state. This creates a stable Frozen Core in equilibrium with the surrounding field, establishing that the interior of a black hole is not a point of infinite density, but a finite region of maximum displacement and absolute zero temperature.


\subsection{Derivation of the Satiation Identity}
Using $G_{eff} = R^3 c^2 / 6M$ and the Schwarzschild boundary condition $R = 2GM/c^2$:
\begin{equation}
R = \frac{2 \left( \frac{R^3 c^2}{6M} \right) M}{c^2} \implies R = \frac{R^3}{3}
\end{equation}
This identity shows that the Frozen Core is a geometric constant of the field's satiation limit. Any constituent of the field that achieves the density required for maximum displacement transitions into this dark, non-vibrational state.

\section{Conclusion}
Dark matter is the description of matter that has reached maximum density, thereby displacing Planck's Field to its absolute limit. These Frozen Cores are scale-invariant, non-vibrational, and thermodynamically at absolute zero. This framework identifies the neutrino and the black hole as identical states of matter existing at disparate scales.

\appendix
\section{The Baryonic-to-Dark Transition Boundary}
The phase transition from baryonic (elastic) matter to dark (frozen) matter is governed by the saturation of the medium's displacement capacity.

\subsection*{A1: Baryonic Phase ($\rho_m < \rho_{sat}$)}
In this phase, the field retains elastic degrees of freedom. The speed of light $v_{\gamma}$ and thermal vibrations $k_B T$ are supported by the medium's ability to oscillate.

\subsection*{A2: Dark Phase ($\rho_m \geq \rho_{sat}$)}
When matter reaches the critical density $\rho_{sat} \approx 5.358 \times 10^{25}$ kg/m$^3$, the field reaches maximum displacement. At this boundary:
\begin{equation}
\Delta \text{Vibration} \to 0 \implies T \to 0 \text{ K}
\end{equation}
The constituent is effectively locked, becoming a non-interactive gravitational node. This transition defines the physical boundary of Dark Matter.

\begin{thebibliography}{9}
\bibitem{Skavberg2026a} Skavberg, R. (2026). \textit{A Mechanical Route to G: Matching Dark-Matter Field Pressure to the Einstein-Hilbert Coupling}. Zenodo. DOI: 10.5281/zenodo.18393278
\bibitem{Skavberg2026b} Skavberg, R. (2026). \textit{A Mechanical Extension of General Relativity: The Singularity in Equilibrium}. Zenodo. DOI: 10.5281/zenodo.18462909
\end{thebibliography}

\end{document}